\documentclass{article}

\usepackage{amsmath,amssymb}
\usepackage{xurl}
\usepackage[hidelinks]{hyperref}
\setlength{\emergencystretch}{2em}

% Shared commands for notes in docs/paper-gaps/.

\DeclareUnicodeCharacter{2080}{\ifmmode{}_{0}\else\textsubscript{0}\fi}
\DeclareUnicodeCharacter{2081}{\ifmmode{}_{1}\else\textsubscript{1}\fi}
\DeclareUnicodeCharacter{2082}{\ifmmode{}_{2}\else\textsubscript{2}\fi}
\DeclareUnicodeCharacter{2099}{\ifmmode{}_{n}\else\textsubscript{n}\fi}

\newcommand{\Lean}{\textsc{Lean}}
\ExplSyntaxOn
\NewDocumentCommand{\leanid}{m}{
  \ifmmode
    \textsf{\detokenize{#1}}
  \else
    \group_begin:
    \urlstyle{sf}
    \tl_set:Nn \l_tmpa_tl {#1}
    \tl_replace_all:Nnn \l_tmpa_tl {\_} {_}
    \exp_args:NV \path \l_tmpa_tl
    \group_end:
  \fi
}
\ExplSyntaxOff
\providecommand{\tr}{\operatorname{tr}}
\newcommand{\tnleanIssueBase}{https://github.com/LionSR/TNLean/issues/}
\newcommand{\tnleanPrBase}{https://github.com/LionSR/TNLean/pull/}
\newcommand{\tnleanDocBase}{https://sirui-lu.com/TNLean/docs/}
\newcommand{\ghissue}[1]{\href{\tnleanIssueBase#1}{GitHub issue \##1}}
\newcommand{\ghpr}[1]{\href{\tnleanPrBase#1}{PR \##1}}
\newcommand{\doclink}[1]{\href{\tnleanDocBase#1.html}{\path{#1}}}

% Dirac notation and the identity operator, as in blueprint/src/macros/common.tex.
% \providecommand keeps note-local definitions authoritative.
\usepackage{dsfont}
\providecommand{\ket}[1]{|#1\rangle}
\providecommand{\bra}[1]{\langle #1|}
\providecommand{\braket}[2]{\langle #1 | #2 \rangle}
\providecommand{\ketbra}[2]{|#1\rangle\!\langle #2|}
\providecommand{\Id}{\mathds{1}}
\providecommand{\one}{\mathds{1}}
\providecommand{\C}{\mathbb{C}}
\providecommand{\MN}[1]{M_{#1}(\C)}
\providecommand{\MD}{M_D(\C)}

% External referencing for paper sources.  The build helper writes sanitized
% auxiliary files under build/paper-aux/ and excludes duplicated source labels.
\usepackage{xr}

% Import a generated paper auxiliary file with a caller-supplied label prefix.
% The candidate paths support builds from docs/paper-gaps/, the repository root,
% and build/paper-aux/ itself.
\newcommand{\importpaperaux}[2]{%
  \IfFileExists{../../build/paper-aux/#2.aux}{%
    \externaldocument[#1]{../../build/paper-aux/#2}%
  }{%
    \IfFileExists{build/paper-aux/#2.aux}{%
      \externaldocument[#1]{build/paper-aux/#2}%
    }{%
      \IfFileExists{#2.aux}{%
        \externaldocument[#1]{#2}%
      }{}%
    }%
  }%
}

% General external reference macros
\newcommand{\paperref}[2]{\ref{#1-#2}}
\newcommand{\papereqref}[2]{(\ref{#1-#2})}

% CPSV16 source (1606.00608 / MPDO-22-12-17-2.tex) external document and shortcuts
\importpaperaux{cpsv16-}{1606.00608/MPDO-22-12-17-2-tnlean}
\newcommand{\cpsvref}[1]{\paperref{cpsv16}{#1}}
\newcommand{\cpsveqref}[1]{\papereqref{cpsv16}{#1}}

% Verdict marker: \gapnote{<kind>}{<status>}. Kind is the policy
% classification (clarification, local-correction, scope-restriction,
% unfaithful, false-source, open-gap); status is open, wip (elimination
% underway), resolved, or historical. Machine-read by the site index;
% prints a verdict line.
\newcommand{\gapnote}[2]{\par\noindent\textbf{Verdict:} #1 (#2).\par}


\title{The Vertical Canonical Form Retains Only the Nonzero Physical Sectors}
\author{}
\date{2026-07-13}

\begin{document}
\maketitle
\gapnote{local-correction}{resolved}

\section{The dimension allowed by canonical form}

The canonical-form construction of
Ref.~\cite{Cirac2016MPDO_arXiv} decomposes a tensor into nonzero irreducible
blocks of dimensions $D_k$. It permits zero blocks in the original
$D$-dimensional space and therefore states, at source lines~214--225, that
\begin{align}
    \sum_k D_k
        &\leq D.
    \label{eq:cpgsv17_vertical_isometry_zero_sector_dimension_bound}
\end{align}
Proposition~4.13 applies this construction to the vertically viewed tensor
$\widetilde M$ and obtains, at source lines~943--951 and~1863--1870,
\begin{align}
    U\widetilde M U^\dagger
        &=
        \bigoplus_\alpha\mu_\alpha\otimes M_\alpha.
    \label{eq:cpgsv17_vertical_isometry_zero_sector_vertical_form}
\end{align}
Here each $\mu_\alpha$ is positive and diagonal, and the tensors
$M_\alpha$ form a basis of normal tensors
\cite[Proposition~4.13]{Cirac2016MPDO_arXiv}.

Let $d$ be the original physical dimension, and let $s$ be the dimension of
the direct sum in
Eq.~\ref{eq:cpgsv17_vertical_isometry_zero_sector_vertical_form}. Applied in
the vertical direction,
Eq.~\ref{eq:cpgsv17_vertical_isometry_zero_sector_dimension_bound} gives
$s\leq d$. In the displayed orientation, the map
$U:\mathbb C^d\to\mathbb C^s$ sends the physical space onto the retained
nonzero-sector space. Hence
\begin{align}
    UU^\dagger
        &=I_s,
    \label{eq:cpgsv17_vertical_isometry_zero_sector_coisometry}\\
    U^\dagger U
        &=P,
    \label{eq:cpgsv17_vertical_isometry_zero_sector_support_projection}
\end{align}
where $P$ is the support projection of the retained sectors. The identity
$U^\dagger U=I_d$ holds only when the zero complement is absent. It is neither
a hypothesis nor a conclusion of Proposition~4.13 in
Ref.~\cite{Cirac2016MPDO_arXiv}.

The source does not print a separate reverse identity. Its preceding statement
that $\widetilde M$ is in canonical form requires the discarded complement to
be a zero sector, which is recorded by
\begin{align}
    \widetilde M
        &=
        U^\dagger
        \left(\bigoplus_\alpha\mu_\alpha\otimes M_\alpha\right)U.
    \label{eq:cpgsv17_vertical_isometry_zero_sector_reconstruction}
\end{align}
Without
Eq.~\ref{eq:cpgsv17_vertical_isometry_zero_sector_reconstruction},
Eq.~\ref{eq:cpgsv17_vertical_isometry_zero_sector_vertical_form} is only a
compression; it does not exclude nonzero entries in the discarded corner or
the off-diagonal corners.

\section{A zero-sector example}

Take horizontal bond dimension one and physical dimension two. Define
\begin{align}
    M^{00}
        &=1,
    \label{eq:cpgsv17_vertical_isometry_zero_sector_rank_one_nonzero_entry}\\
    M^{01}
        &=M^{10}
        =M^{11}
        =0.
    \label{eq:cpgsv17_vertical_isometry_zero_sector_rank_one_zero_entries}
\end{align}
This is a one-block horizontal normal tensor. At every positive length $N$, it
generates the positive operator
\begin{align}
    \rho^{(N)}(M)
        &=
        \ket{0\cdots0}\bra{0\cdots0}.
    \label{eq:cpgsv17_vertical_isometry_zero_sector_rank_one_operator}
\end{align}
The vertically viewed tensor has the single letter
\begin{align}
    \widetilde M
        &=
        \begin{pmatrix}
            1&0\\
            0&0
        \end{pmatrix}.
    \label{eq:cpgsv17_vertical_isometry_zero_sector_vertical_letter}
\end{align}
Its nonzero canonical part is the one-dimensional tensor whose letter is $1$.
Equation~\ref{eq:cpgsv17_vertical_isometry_zero_sector_vertical_form} holds
with
\begin{align}
    U
        &=
        \begin{pmatrix}
            1&0
        \end{pmatrix},
    \label{eq:cpgsv17_vertical_isometry_zero_sector_example_map}\\
    UU^\dagger
        &=1,
    \label{eq:cpgsv17_vertical_isometry_zero_sector_example_coisometry}\\
    U^\dagger U
        &=
        \begin{pmatrix}
            1&0\\
            0&0
        \end{pmatrix}.
    \label{eq:cpgsv17_vertical_isometry_zero_sector_example_support}
\end{align}

A square unitary description can retain an explicit zero block $[0]$. The
earlier formal predicate did not include such a block: its right-hand side
contained only strictly positive weights multiplying normal tensors. Within
that positive-block-only assembly, a unitary preserves the two-dimensional
space, but the second one-dimensional corner of
Eq.~\ref{eq:cpgsv17_vertical_isometry_zero_sector_vertical_letter} necessarily
has weight zero. Treating the full matrix as one block does not produce a
normal tensor. Thus a two-sided unitary combined with the positive-block-only
assembly would make Proposition~4.13 false even in this elementary example.

\section{Trace loss under transported sector maps}

Appendix~C.4 of Ref.~\cite{Cirac2016MPDO_arXiv} transports the physical
refinement and coarse-graining maps to the retained one-site and two-site
vertical coordinates at source lines~1955--1980. Let $U_1$ and $U_2$ be the
corresponding coisometries, and define
\begin{align}
    P_1
        &=U_1^\dagger U_1,
    \label{eq:cpgsv17_vertical_isometry_zero_sector_one_site_projection}\\
    P_2
        &=U_2^\dagger U_2.
    \label{eq:cpgsv17_vertical_isometry_zero_sector_two_site_projection}
\end{align}
The normalized diagonal embedding and the sectorwise partial trace preserve
the sum of sector traces, and the physical maps preserve trace. Let $Z_T\geq0$
be the matrix immediately before the final $U_2$ compression in the
transported refinement map $\widetilde T$. Then
\begin{align}
    \tr_{\mathrm{sec}}(\widetilde T(X))
        &=
        \tr(U_2Z_TU_2^\dagger),
    \label{eq:cpgsv17_vertical_isometry_zero_sector_refinement_trace_compression}\\
        &=
        \tr(P_2Z_T),
    \label{eq:cpgsv17_vertical_isometry_zero_sector_refinement_trace_projection}\\
        &\leq
        \tr(Z_T)
        =
        \tr_{\mathrm{sec}}(X).
    \label{eq:cpgsv17_vertical_isometry_zero_sector_refinement_trace_loss}
\end{align}
For positive $Z_T$, equality in
Eq.~\ref{eq:cpgsv17_vertical_isometry_zero_sector_refinement_trace_loss}
holds if and only if
\begin{align}
    P_2Z_T
        &=Z_T.
    \label{eq:cpgsv17_vertical_isometry_zero_sector_refinement_support}
\end{align}

For the transported coarse-graining map $\widetilde S$, let $Z_S\geq0$ denote
the matrix before the final $U_1$ compression. The same argument gives
\begin{align}
    \tr_{\mathrm{sec}}(\widetilde S(Y))
        &\leq
        \tr_{\mathrm{sec}}(Y),
    \label{eq:cpgsv17_vertical_isometry_zero_sector_coarse_trace_loss}\\
    \tr_{\mathrm{sec}}(\widetilde S(Y))
        &=
        \tr_{\mathrm{sec}}(Y)
        \quad\Longleftrightarrow\quad
        P_1Z_S=Z_S.
    \label{eq:cpgsv17_vertical_isometry_zero_sector_coarse_trace_equality}
\end{align}

Consequently, trace preservation of the physical maps does not imply that the
transported maps preserve trace on arbitrary elements of the full finite
products. That conclusion would require
\begin{align}
    P_2Z_T
        &=Z_T,
    \label{eq:cpgsv17_vertical_isometry_zero_sector_transport_refinement_image}\\
    P_1Z_S
        &=Z_S.
    \label{eq:cpgsv17_vertical_isometry_zero_sector_transport_coarse_image}
\end{align}
The identities displayed in Appendix~C.4 of
Ref.~\cite{Cirac2016MPDO_arXiv} establish the required action on the weighted
sector operators arising from the canonical forms. By themselves, they do not
establish
Eqs.~\ref{eq:cpgsv17_vertical_isometry_zero_sector_transport_refinement_image}
and~\ref{eq:cpgsv17_vertical_isometry_zero_sector_transport_coarse_image}
for every element of the finite products.

For every virtual bond matrix $X$, the source-generated families now satisfy
these support identities. The normalized embedding produces the contracted
weighted direct sum; the two reconstruction identities identify the
precompression matrix with $U_i^\dagger B_i(X)U_i$; and
$U_iU_i^\dagger=I$ gives the support equation. Cyclicity of trace then gives
trace equality on these families. No global trace-preserving extension,
complete-positivity assertion for the finite products, or inverse relation is
claimed.

\section{Ambient completion in the converse construction}

Definition~4.1 of Ref.~\cite{Cirac2016MPDO_arXiv} requires trace-preserving
completely positive maps on the full one-site and two-site physical matrix
algebras. In the converse implication of Theorem~4.14, Appendix~C.4 first
defines, at source lines~2065--2085, the retained-sector compositions
\begin{align}
    T_0
        &=\widetilde R_2\widetilde T R_1,
    \label{eq:cpgsv17_vertical_isometry_zero_sector_active_refinement_composite}\\
    S_0
        &=\widetilde R_1\widetilde S R_2.
    \label{eq:cpgsv17_vertical_isometry_zero_sector_active_coarse_composite}
\end{align}
The maps $R_1$ and $R_2$ begin by compressing with the one-site and two-site
coisometries. After the canonical relabelling of the two-site indices, with
$P_i=U_i^\dagger U_i$, one obtains
\begin{align}
    \tr(T_0(X))
        &=\tr(P_1X),
    \label{eq:cpgsv17_vertical_isometry_zero_sector_active_refinement_trace}\\
    \tr(S_0(Y))
        &=\tr(P_2Y).
    \label{eq:cpgsv17_vertical_isometry_zero_sector_active_coarse_trace}
\end{align}
These maps preserve trace on matrices supported in the retained sectors, but
not on arbitrary ambient matrices when $P_i\neq I$.

Define the complementary projections by
\begin{align}
    Q_i
        &=I-P_i.
    \label{eq:cpgsv17_vertical_isometry_zero_sector_complement_projection}
\end{align}
When $d>0$, choose density matrices $\tau_1$ and $\tau_2$ on the one-site and
two-site physical spaces. Define the measure-and-prepare completions
\begin{align}
    T(X)
        &=
        T_0(X)+\tr(Q_1XQ_1)\tau_2,
    \label{eq:cpgsv17_vertical_isometry_zero_sector_completed_refinement_channel}\\
    S(Y)
        &=
        S_0(Y)+\tr(Q_2YQ_2)\tau_1.
    \label{eq:cpgsv17_vertical_isometry_zero_sector_completed_coarse_channel}
\end{align}
Both maps are completely positive. Since each $Q_i$ is a projection and
$\tr(\tau_i)=1$,
Eqs.~\ref{eq:cpgsv17_vertical_isometry_zero_sector_active_refinement_trace}
and~\ref{eq:cpgsv17_vertical_isometry_zero_sector_active_coarse_trace} give
\begin{align}
    \tr(T(X))
        &=\tr(X),
    \label{eq:cpgsv17_vertical_isometry_zero_sector_completed_refinement_trace}\\
    \tr(S(Y))
        &=\tr(Y).
    \label{eq:cpgsv17_vertical_isometry_zero_sector_completed_coarse_trace}
\end{align}
When $d=0$, both ambient physical matrix algebras contain only zero. The
uncompleted maps are then already completely positive and trace preserving,
and no density matrix is needed.

The source does not specify $\tau_i$ or any other trace-restoring action on the
discarded complements. This freedom does not affect its conclusion. For every
virtual matrix $Z$, reconstruction gives
\begin{align}
    P_1M_1(Z)P_1
        &=M_1(Z),
    \label{eq:cpgsv17_vertical_isometry_zero_sector_one_site_source_support}\\
    P_2M_2(Z)P_2
        &=M_2(Z).
    \label{eq:cpgsv17_vertical_isometry_zero_sector_two_site_source_support}
\end{align}
The added terms in
Eqs.~\ref{eq:cpgsv17_vertical_isometry_zero_sector_completed_refinement_channel}
and~\ref{eq:cpgsv17_vertical_isometry_zero_sector_completed_coarse_channel}
therefore vanish on the source-generated family to which source
lines~2074 and~2085 are applied.

This completion differs from the transported maps of the preceding section.
In the forward implication at source lines~1955--1980, full physical channels
are hypotheses and are compressed to retained-sector coordinates; their
possible trace loss is governed by
Eqs.~\ref{eq:cpgsv17_vertical_isometry_zero_sector_transport_refinement_image}
and~\ref{eq:cpgsv17_vertical_isometry_zero_sector_transport_coarse_image}. In
the converse implication at source lines~2065--2085, sector maps are used to
construct full physical channels, and
Eqs.~\ref{eq:cpgsv17_vertical_isometry_zero_sector_completed_refinement_channel}
and~\ref{eq:cpgsv17_vertical_isometry_zero_sector_completed_coarse_channel}
supply the otherwise unspecified action on the orthogonal complements.

\section{Formal correction}

The vertical canonical-form predicate now requires
Eq.~\ref{eq:cpgsv17_vertical_isometry_zero_sector_coisometry} in the
orientation of
Eq.~\ref{eq:cpgsv17_vertical_isometry_zero_sector_vertical_form}, together
with exact reconstruction from
Eq.~\ref{eq:cpgsv17_vertical_isometry_zero_sector_reconstruction}. This
conclusion is compatible with
Eq.~\ref{eq:cpgsv17_vertical_isometry_zero_sector_dimension_bound} and is
faithful to Proposition~4.13 of Ref.~\cite{Cirac2016MPDO_arXiv}: the proposition
prints the compression, while its preceding canonical-form assertion supplies
reconstruction. The retained part is represented rectangularly, and the
omitted part must be zero. No full-support hypothesis is added.

Literal Proposition~4.13 now uses this condition in
\path{TNLean/MPS/MPDO/CPSVVerticalCanonicalForm.lean}.\footnote{Formal
declaration: \leanid{MPOTensor.verticalCF\_of\_cpsvCanonicalForm}.}
The theorem under the normalized BNT-refined horizontal form remains available
separately in
\path{TNLean/MPS/MPDO/VerticalCanonicalForm.lean}.\footnote{Formal declaration:
\leanid{MPOTensor.verticalCF\_of\_horizontalCF}.}
The vertical phase-class grouping and normalization preceding both coisometry
constructions are recorded in
\path{cpgsv17_vertical_cf_grouping.tex}. The coisometry correction is tracked
in \href{https://github.com/LionSR/TNLean/issues/3869}{GitHub issue \#3869}.

The source-generated support identities and their trace consequences are
proved in
\path{TNLean/MPS/MPDO/VerticalSectorImagePreservation.lean}.\footnote{The image
identities are
\leanid{MPOTensor.precompressionVerticalSectorT_image_preservation} and
\leanid{MPOTensor.precompressionVerticalSectorS_image_preservation}; their
trace consequences are
\leanid{MPOTensor.transportedVerticalSectorT_trace_of_source_generated} and
\leanid{MPOTensor.transportedVerticalSectorS_trace_of_source_generated}.}

The general measure-and-prepare completion is formalized in
\path{TNLean/Channel/SupportCompletion.lean}.\footnote{Formal declaration:
\leanid{Matrix.supportCompletion\_isKrausCPTP}.}
Under supplied unitary conjugacies between the paired sectors, the active
compositions and completed physical channels are formalized in
\path{TNLean/MPS/MPDO/BNTAlgebraTensorClauseConditionalPhysicalMaps.lean}.%
\footnote{Formal declaration:
\leanid{MPOTensor.BNTAlgebraTensorClause.TwoSiteExactSectorGauge.%
UnitarySectorConjugacy.isRFPViaTS}.}
The existence of these unitary conjugacies is treated in
\path{cpsv16_two_site_sector_unitary_gauge_gap.tex}. The complementary terms
are choices on discarded subspaces, not additional claims about Appendix~C.4
of Ref.~\cite{Cirac2016MPDO_arXiv}. Their action on the source-generated family
is fixed by
Eqs.~\ref{eq:cpgsv17_vertical_isometry_zero_sector_one_site_source_support}
and~\ref{eq:cpgsv17_vertical_isometry_zero_sector_two_site_source_support}.

\paragraph{Verdict.}
The earlier two-sided-unitary condition was absent from the source and failed
when the vertically viewed tensor contained a zero sector. Replacing it with
the coisometry identity $UU^\dagger=I$, while retaining exact reconstruction
from the nonzero sectors, corrects the conclusion without adding a full-support
hypothesis. Literal Proposition~4.13 and the separate stronger horizontal
theorem now prove this corrected conclusion.

This completion does not make the forward transported maps trace preserving on
every element of the finite products in Appendix~C.4; the support conditions
in
Eqs.~\ref{eq:cpgsv17_vertical_isometry_zero_sector_transport_refinement_image}
and~\ref{eq:cpgsv17_vertical_isometry_zero_sector_transport_coarse_image}
remain necessary. It also does not choose a preferred action on the discarded
physical complements.

For the converse implication, the measure-and-prepare terms in
Eqs.~\ref{eq:cpgsv17_vertical_isometry_zero_sector_completed_refinement_channel}
and~\ref{eq:cpgsv17_vertical_isometry_zero_sector_completed_coarse_channel}
produce global physical channels with the same action on the source-generated
family. Within this note, the construction is conditional on the unitary
sector conjugacies used in
Eqs.~\ref{eq:cpgsv17_vertical_isometry_zero_sector_active_refinement_composite}
and~\ref{eq:cpgsv17_vertical_isometry_zero_sector_active_coarse_composite}.
The later mixed one-site/two-site prefix argument proves these conjugacies from
the tensor-attached BNT clause, as documented in
\path{cpsv16_two_site_sector_unitary_gauge_gap.tex}. Consequently, the
zero-sector completion is used unconditionally by
\leanid{MPOTensor.BNTAlgebraTensorClause.isRFPViaTS}, closing
Theorem~4.14(ii)$\Rightarrow$(i) under the standing canonical-form and MPDO
hypotheses of Ref.~\cite{Cirac2016MPDO_arXiv}.

\bibliographystyle{alpha}
\bibliography{references}

\end{document}
